\chapter{Floorplanning and Chip Setup}
\label{ch:floorplan}

\glance{Floorplanning turns a linked netlist into a physical canvas: die and
core areas, standard-cell rows, routing tracks, well taps / endcaps, optionally
an I/O pad ring or bump array, and power-intent domains from UPF. Modules:
\pkg{src/ifp}, \pkg{src/tap}, \pkg{src/pad}, \pkg{src/upf}.}

\section{Initialize floorplan: \pkg{src/ifp}}

\id{ifp::InitFloorplan} creates the die/core geometry and the placement rows.
Two mutually exclusive sizing styles:

\begin{itemize}
  \item \textbf{Utilization:} \cmd{-utilization}, \cmd{-aspect\_ratio},
        \cmd{-core\_space} --- die/core sizes are derived from design area.
  \item \textbf{Explicit:} \cmd{-die\_area} / \cmd{-core\_area} as rectangles
        or rectilinear polygons (L/T shapes).
\end{itemize}

\begin{lstlisting}[style=ortcl]
initialize_floorplan -utilization 40 -aspect_ratio 1.0 \
    -core_space 2.0 -site FreePDK45_38x28_10R_NP_162NW_34O
make_tracks            ;# create routing tracks on each metal layer
\end{lstlisting}

\cmd{-site} selects the single-height site whose rows are built; additional /
flip sites support hybrid-row and multi-height designs. Rows are written into
odb as \id{dbRow} objects; \cmd{make\_tracks} creates the preferred-direction
routing tracks from the LEF technology.

\section{Tapcells and endcaps: \pkg{src/tap}}

\id{tap::Tapcell} inserts well-tap and endcap cells into the rows for latch-up
prevention and DRC. The main command is \cmd{tapcell}:

\begin{lstlisting}[style=ortcl]
tapcell -distance 120 \
        -tapcell_master TAPCELL_X1 \
        -endcap_master  ENDCAP_XL
\end{lstlisting}

\cmd{-distance} sets the checkerboard spacing (microns). Extra masters
(\cmd{-cnrcap\_*}, \cmd{-tap\_nwin*}, \ldots) handle corners and macro
boundaries. \cmd{cut\_rows} trims rows around macros with a halo.

\section{I/O pads and bumps: \pkg{src/pad} (ICeWall)}

\pkg{src/pad} builds chip-level I/O: a pad ring around the die and/or a bump
array for flip-chip. Representative commands:

\begin{center}
\small
\begin{tabularx}{\linewidth}{@{}lX@{}}
\toprule
\textbf{Command} & \textbf{Purpose} \\
\midrule
\cmd{make\_io\_sites} / \cmd{place\_pad} / \cmd{place\_pads} & build / place the padring \\
\cmd{make\_io\_bump\_array} / \cmd{place\_bondpad} & bump / bond-pad arrays \\
\cmd{place\_io\_terminals} & place terminals on padcell pins \\
\cmd{connect\_by\_abutment} & abutting power/ground pad connections \\
\bottomrule
\end{tabularx}
\end{center}

\section{Power intent: \pkg{src/upf}}

\pkg{src/upf} reads Unified Power Format into odb power-domain objects:

\begin{lstlisting}[style=ortcl]
read_upf -file design.upf
write_upf design_out.upf
\end{lstlisting}

UPF describes power domains, supply nets, isolation / level-shifter cells, and
related constraints that later stages (placement, PDN) must respect.

\section{Summary}

\begin{table}[h]
\centering
\small
\begin{tabularx}{\linewidth}{@{}lllX@{}}
\toprule
\textbf{Step} & \textbf{Module} & \textbf{Command} & \textbf{Writes to DB} \\
\midrule
Die/core/rows & \pkg{ifp} & \cmd{initialize\_floorplan} & die/core, \id{dbRow} \\
Tracks        & \pkg{ifp} & \cmd{make\_tracks} & routing tracks \\
Taps/endcaps  & \pkg{tap} & \cmd{tapcell} & tap/endcap \id{dbInst} \\
I/O ring      & \pkg{pad} & \cmd{place\_pads}, \ldots & pad/bump instances \\
Power intent  & \pkg{upf} & \cmd{read\_upf} & power domains \\
\bottomrule
\end{tabularx}
\end{table}
